\begin{derivation}[从\eqnrefs{eq:sm-top-bw-amplitude-dp94}到\eqnrefs{eq:sm-top-helicity-fractions-dp94}]
在顶夸克静止系取
\begin{equation*}
 p_t^\mu=(m_tc,0,0,0),
 \qquad
 k^\mu=\left(E_W/c,0,0,+|\mathbf k|\right),
 \qquad
 p_b^\mu=\left(E_b/c,0,0,-|\mathbf k|\right),
\end{equation*}
且 $m_b=0$。二体运动学给
\begin{equation*}
 E_W=\frac{m_t^2+m_W^2}{2m_t}c^2,
 \qquad
 E_b=|\mathbf k|c=\frac{m_t^2-m_W^2}{2m_t}c^2.
\end{equation*}
对 $W$ 取极化矢量
\begin{align*}
 \varepsilon_0^\mu
 &=\frac1{m_Wc}\left(|\mathbf k|,0,0,E_W/c\right),\\
 \varepsilon_\pm^\mu
 &=\frac1{\sqrt2}(0,\mp1,-\ii,0),
\end{align*}
满足 $k\cdot\varepsilon_\lambda=0$ 与 $\varepsilon_\lambda^*\cdot\varepsilon_\lambda=-1$。

固定 $W$ 极化后，对末态 $b$ 自旋求和并对初态顶夸克自旋平均。利用
\(\sum_su_s(p)\bar u_s(p)=\slashed p+mc\)，有
\begin{equation*}
 \overline{|\mathcal M_\lambda|^2}
 =\frac{g^2|V_{tb}|^2}{4}
 \varepsilon_{\mu}^{(\lambda)*}\varepsilon_\nu^{(\lambda)}
 L^{\mu\nu},
\qquad
 L^{\mu\nu}\equiv
 \Tr\!\left[
 \slashed p_b\gamma^\mu P_L(\slashed p_t+m_tc)
 \gamma^\nu P_L
 \right].
\end{equation*}
先把这个迹本身算清楚。由于
\(P_L\gamma^\nu=\gamma^\nu P_R\) 且 $P_RP_L=0$，质量项满足
\begin{equation*}
 m_tc\,\Tr(\slashed p_b\gamma^\mu P_L\gamma^\nu P_L)=0.
\end{equation*}
所以只剩四动量项。展开两个手征投影，含单个 $\gamma^5$ 的部分用
\(\Tr(\gamma^\alpha\gamma^\mu\gamma^\beta\gamma^\nu\gamma^5)
=-4\ii\epsilon^{\alpha\mu\beta\nu}\)，不含 $\gamma^5$ 的部分用四伽马迹，得到
\begin{equation*}
\boxed{
L^{\mu\nu}
=2\left[
 p_b^\mu p_t^\nu+p_b^\nu p_t^\mu
 -\eta^{\mu\nu}(p_b\cdot p_t)
 -\ii\epsilon^{\mu\nu\alpha\beta}p_{b\alpha}p_{t\beta}
 \right].}
\end{equation*}
这一步把“$V-A$”的两种数学作用分开了：前三项来自矢量迹，最后的反对称项来自 $\gamma^5$，它正负责区分两个横向螺旋度。

在顶夸克静止系中记
\begin{equation*}
 Q\equiv p_t\cdot p_b
 =\frac{m_t^2c^2}{2}(1-x_W).
\end{equation*}
代入上面的 $p_t,p_b$ 后，$L^{\mu\nu}$ 的非零部分可直接写成
\begin{equation*}
L^{\mu\nu}
=2Q
\begin{pmatrix}
 1&0&0&-1\\
 0&1&\ii&0\\
 0&-\ii&1&0\\
 -1&0&0&1
\end{pmatrix}^{\mu\nu}.
\end{equation*}
现在三个螺旋度无需再借助口头角动量论证，而可以逐个收缩。

对正螺旋度，协变偏振分量为
\begin{equation*}
\varepsilon_{\mu}^{(+)}
=\frac1{\sqrt2}(0,1,\ii,0).
\end{equation*}
横向 $x$--$y$ 子块作用在列向量 $(1,\ii)^{\mathsf T}$ 上恰好为零：
\begin{equation*}
\begin{pmatrix}1&\ii\\-\ii&1\end{pmatrix}
\begin{pmatrix}1\\\ii\end{pmatrix}
=\begin{pmatrix}0\\0\end{pmatrix}.
\end{equation*}
因此
\begin{equation*}
\boxed{\overline{|\mathcal M_+|^2}=0}.
\end{equation*}
这里的零来自手征张量中对称部分与 $\gamma^5$ 反对称部分的精确相消；这就是 $m_b=0$ 时右手 $W$ 被禁止的代数形式。

对负螺旋度，
\(\varepsilon_{\mu}^{(-)}=(0,-1,\ii,0)/\sqrt2\)，于是
\begin{align*}
\varepsilon_{\mu}^{(-)*}L^{\mu\nu}\varepsilon_{\nu}^{(-)}
&=4Q
=2m_t^2c^2(1-x_W).
\end{align*}
按正文采用的约化振幅归一化，这给出
\begin{equation*}
 \boxed{
 \overline{|\mathcal M_-|^2}
 =\frac{g^2|V_{tb}|^2}{2}
 m_t^2(1-x_W)}.
\end{equation*}

纵向偏振只有 $0,z$ 分量。令
\begin{equation*}
 a\equiv\frac{|\mathbf k|}{m_Wc},
 \qquad
 b\equiv\frac{E_W}{m_Wc^2},
 \qquad
 \varepsilon_{\mu}^{(0)}=(a,0,0,-b).
\end{equation*}
$0$--$z$ 子块同样可以直接收缩：
\begin{align*}
\varepsilon_{\mu}^{(0)}L^{\mu\nu}\varepsilon_{\nu}^{(0)}
&=2Q(a+b)^2.
\end{align*}
而二体运动学给
\begin{align*}
 a+b
&=\frac{|\mathbf k|c+E_W}{m_Wc^2}\\
&=\frac{m_tc^2(1-x_W)/2+m_tc^2(1+x_W)/2}{m_Wc^2}\\
&=\frac{m_t}{m_W}.
\end{align*}
因此
\begin{equation*}
\varepsilon_{\mu}^{(0)*}L^{\mu\nu}\varepsilon_{\nu}^{(0)}
=\frac{2Q}{x_W},
\end{equation*}
从而
\begin{equation*}
 \boxed{
 \overline{|\mathcal M_0|^2}
 =\frac{g^2|V_{tb}|^2}{4}
 \frac{m_t^4}{m_W^2}(1-x_W)}.
\end{equation*}
纵向增强的来源现在是显式的：$\varepsilon_0^\mu$ 自身含 $1/m_W$，而 $a+b=m_t/m_W$，所以当 $m_t\gg m_W$ 时纵向通道相对于横向通道增强。
因此
\begin{equation*}
 \frac{\overline{|\mathcal M_-|^2}}
 {\overline{|\mathcal M_0|^2}}
 =2\frac{m_W^2}{m_t^2}=2x_W,
\end{equation*}
而 $\overline{|\mathcal M_+|^2}=0$。用
\begin{equation*}
F_\lambda
=\frac{\overline{|\mathcal M_\lambda|^2}}
{\overline{|\mathcal M_0|^2}+\overline{|\mathcal M_-|^2}+\overline{|\mathcal M_+|^2}}
\end{equation*}
逐项除以 $\overline{|\mathcal M_0|^2}$，得到
\begin{equation*}
 F_0=\frac1{1+2x_W},
 \qquad
 F_-=\frac{2x_W}{1+2x_W},
 \qquad
 F_+=0.
\end{equation*}

总自旋平均振幅平方为
\begin{equation*}
 \overline{|\mathcal M|^2}
 =\frac{g^2|V_{tb}|^2}{4}
 \frac{m_t^4}{m_W^2}(1-x_W)(1+2x_W).
\end{equation*}
二体衰变公式
\begin{equation*}
 \Gamma
 =\frac{|\mathbf k|}{8\pi\hbar m_t^2c}
 \overline{|\mathcal M|^2}
\end{equation*}
与 $|\mathbf k|=m_tc(1-x_W)/2$ 合并，再用
\begin{equation*}
 \frac{G_F^{(E)}}{\sqrt2}=\frac{g^2}{8(m_Wc^2)^2},
\end{equation*}
即得到正文\eqnrefs{eq:sm-top-width-dp94}。
\end{derivation}
